% ============================================================================
% 第 2 章 相关技术概述
% 预估页数：3-4 页
% 写作总要点：
%   - 介绍系统所用的关键技术，每项 200-400 字
%   - 重点是为后续章节（特别是 Agent 章节）做理论铺垫
%   - 不要写成技术说明书，要点出"为何选用 + 在本系统中的作用"
% ============================================================================
\chapter{相关技术概述}

% ----------------------------------------------------------------------------
\section{前端开发技术}

\subsection{Vue 3 与 Composition API}
% 写作要点：
%   - 渐进式框架、响应式系统的优势
%   - Composition API 相对 Options API 在大型项目的逻辑复用优势
%   - 引用 Vue 官方文档 \cite{vuejs2024}

Vue.js 是渐进式 JavaScript 前端框架，适合以组件方式组织中后台页面\cite{vuejs2024}。本系统采用 Vue 3.4.21，其基于 Proxy 的响应式机制更适合处理训练页中不断变化的对话列表、计时状态和表单数据。相比 Options API，Composition API 可以把同一业务关注点集中到一个逻辑单元中，例如将 Agent 状态、SSE 对话流、训练提交和病例信息拆成独立复合函数，避免大型页面的状态散落在多个选项中。系统同时使用 Vue Router、Pinia 和 Vite 完成路由、状态管理与开发构建。

\subsection{Element Plus 与 ECharts}
% 写作要点：
%   - Element Plus 提供企业级 UI 组件
%   - ECharts 用于雷达图（能力评估）、折线图（成绩趋势）、柱状图（知识点掌握）
%   - 引用 \cite{li2018echarts}

Element Plus 是面向 Vue 3 的 UI 组件库，提供表格、表单、对话框、上传和菜单等常用组件。本系统的学生端、教师端和管理员端均基于 Element Plus 2.6.3 构建，并通过 CSS 变量区分不同角色的主题色，减少重复页面开发。

ECharts 用于系统中的学情可视化\cite{echarts2024}。学生端使用雷达图展示问诊、诊断、治疗和沟通四个维度，使用折线图呈现训练成绩变化；教师端则通过柱状图和饼图观察班级知识点错误分布。相比直接展示表格，图表更利于学生复盘和教师发现共性问题。

% ----------------------------------------------------------------------------
\section{后端开发技术}

\subsection{Spring Boot 与 Spring Security}
% 写作要点：
%   - 约定优于配置，自动装配
%   - Spring Security 提供过滤器链、认证授权
%   - 本系统采用 Spring Boot 3.3.6 + Java 17

Spring Boot 在 Spring Framework 基础上提供自动装配和约定式配置，能减少重复工程代码\cite{springboot2024}。本系统采用 Spring Boot 3.3.6 和 Java 17，主要用于提供 REST 接口、SSE 流式对话接口以及与数据库、缓存、对象存储的集成。其 WebFlux 组件被用于调用外部大模型服务，便于处理较长时间的 AI 响应。

Spring Security 负责身份认证和权限控制。本系统在过滤器链中加入 JWT 认证过滤器，从 HTTP 头中读取 Token、验证签名并写入认证上下文；登录、验证码、文件预览等少数路径开放访问，其余接口默认需要认证。该设计能够满足前后端分离系统的基本安全边界。

\subsection{MyBatis-Plus 与 JWT}
% 写作要点：
%   - MyBatis-Plus 在 MyBatis 基础上增强，提供 LambdaQueryWrapper 等
%   - JWT (RFC 7519) \cite{jones2015jwt}：无状态 Token、跨域友好

MyBatis-Plus 在 MyBatis 基础上提供通用 CRUD、类型安全查询、逻辑删除、自动填充和分页插件等能力\cite{mybatisplus2024}。本系统使用 \texttt{LambdaQueryWrapper} 表达查询条件，使用分页插件处理训练记录和批改列表，并通过 \texttt{deleted} 字段保留业务数据的追溯能力。

JWT（JSON Web Token）是一种自包含令牌规范\cite{jones2015jwt}，适合前后端分离和跨域调用。本系统使用 JJWT 完成令牌签发与验证，在 Payload 中保存用户 ID、用户名和过期时间，并通过 Redis 存储刷新令牌；登录失败锁定由 \texttt{LoginAttemptService} 处理，用于降低暴力破解风险。

% ----------------------------------------------------------------------------
\section{大语言模型与阿里百炼平台}
% 写作要点：
%   - 大语言模型的预训练 + 指令微调范式
%   - 通义千问 Qwen 系列模型 \cite{bai2023qwen}
%   - 阿里百炼 DashScope 提供 OpenAI 兼容 API，便于工程集成
%   - 本系统使用 qwen-plus / qwen3.5-flash 作为主模型

大语言模型通常先通过大规模预训练获得通用语言和知识能力，再通过指令微调与偏好对齐改善可控性。通义千问（Qwen）是中文友好的大语言模型系列\cite{bai2023qwen}，对中文对话、代码和多学科知识支持较好，适合作为本系统虚拟病人和训练反馈的基础模型。

阿里云百炼 DashScope 以 SaaS 形式提供通义千问推理服务，并兼容 OpenAI Chat Completions 接口\cite{dashscope2024}。本系统默认使用 \texttt{qwen-plus}，生产环境可切换为响应更快的 \texttt{qwen3.5-flash}。后端通过 \texttt{WebClient} 调用兼容端点，使用 \texttt{Bearer} Token 认证 API Key，并通过流式输出改善长回复等待体验。

% ----------------------------------------------------------------------------
\section{AI Agent 与多智能体协作}
% 写作要点：
%   - Agent 的定义：感知 + 思考 + 行动
%   - ReAct 框架 \cite{yao2023react}：Reasoning + Acting 交错
%   - 多 Agent 协作（multi-agent collaboration）的研究综述 \cite{wang2024llmagent}
%   - 引出本系统的 PatientAgent / DiagnosisAgent / FeedbackAgent 设计

在大语言模型应用中，Agent 通常指以 LLM 为核心，并配合记忆、规划和工具调用的软件实体。ReAct 框架\cite{yao2023react}将推理与行动交替组织，使模型可以根据外部反馈逐步修正任务执行过程。

多智能体协作进一步强调角色拆分。AutoGen\cite{wu2023autogen}和 MetaGPT\cite{hong2024metagpt}均通过不同 Agent 的提示词、工具集和对话协议完成复杂任务。Wang 等\cite{wang2024llmagent}将相关能力概括为感知、记忆、规划和行动，这也为领域专用 Agent 的设计提供了参考。

本系统的 Agent 编排层针对医学问诊场景实现：PatientAgent 在问诊阶段扮演标准化病人，DiagnosisAgent 在提交阶段点评诊断思路，FeedbackAgent 输出治疗反馈与风险提醒。三者由 \texttt{AgentRouter} 按训练阶段选择，避免一个提示词同时负责扮演和评估。

% ----------------------------------------------------------------------------
\section{检索增强生成 RAG}
% 写作要点：
%   - Lewis 等人提出的 RAG 框架 \cite{lewis2020rag}
%   - 解决 LLM 幻觉、知识时效性问题
%   - 经典 RAG = 向量检索 + 重排 + 上下文注入
%   - 本系统采用关键词匹配的轻量级 RAG（不依赖向量库），适合医学领域结构化知识点

大语言模型的知识受训练语料和更新时间限制，生成时也可能给出没有依据的事实。检索增强生成（Retrieval-Augmented Generation，RAG）在生成前先检索外部知识，并把相关内容拼接到提示词中\cite{lewis2020rag}。这种方法不需要重新训练模型，就能把领域知识引入回答过程。

经典 RAG 通常包括向量化、向量检索、重排和上下文拼接，依赖嵌入模型与向量数据库。但本系统的知识点数量不大，主题又按科室、症状和疾病结构化存储，使用复杂向量库的收益有限。因此 \texttt{KnowledgeRAGService} 采用关键词匹配：从用户问题中提取较长关键词，在知识点标题和标签字段上进行模糊查询，取前三条结果拼接为提示词上下文。该方案牺牲了一部分语义召回能力，但部署简单、可解释性强，适合当前规模的问诊实训系统。

% ----------------------------------------------------------------------------
\section{工具调用 Tool Calling / Function Calling}
% 写作要点：
%   - Toolformer \cite{schick2023toolformer}
%   - OpenAI Function Calling 规范：JSON Schema 描述工具
%   - 本系统通过 AiToolService 提供 query\_symptom\_knowledge / get\_diagnosis\_criteria /
%     evaluate\_treatment\_plan 三个工具

工具调用（Tool Calling / Function Calling）用于让大语言模型访问外部服务或结构化数据。Toolformer\cite{schick2023toolformer}讨论了模型学习工具调用的方式，OpenAI 兼容接口则通常使用 \texttt{tools} 和 \texttt{tool\_calls} 字段描述工具名称、参数和调用结果。相比让模型直接生成代码再解析，结构化工具调用更稳定。

本系统的 \texttt{AiToolService} 定义了知识查询、诊断标准获取和治疗方案记录三个医学工具。各工具通过 \texttt{KnowledgePointMapper} 在知识点表上检索并返回结构化结果。考虑到当前系统稳定性优先，配置项 \texttt{enableTools} 默认关闭，后续可在更复杂的训练任务中逐步启用。

% ----------------------------------------------------------------------------
\section{数据库与缓存技术}

\subsection{MySQL 8.0}
% 写作要点：业务数据持久化，支持 JSON 字段（病例附件信息）

本系统采用 MySQL 8.0 作为主业务数据库，使用 InnoDB 存储引擎和 \texttt{utf8mb4} 字符集。病例表中的 \texttt{medical\_images} 字段保存影像附件信息，评分表中的 \texttt{weak\_points}、\texttt{error\_patterns}、\texttt{score\_detail} 等字段保存结构化评分结果。部分动态字段使用 JSON 存储，可减少额外建表。

\subsection{Redis 7}
% 写作要点：会话缓存、分布式锁、Memory 模块用于存储多轮对话上下文

Redis 是内存型键值存储系统\cite{redis2024}。本系统使用 Redis 7 保存验证码、登录失败计数、账号锁定标记和多轮对话上下文。会话记忆服务按训练记录编号保存最近对话文本，并设置 24 小时过期和最多 10 轮的滑动窗口。

\subsection{MinIO 对象存储}
% 写作要点：S3 兼容、医学影像与病例附件存储

MinIO 是兼容 Amazon S3 API 的对象存储服务\cite{minio2024}，用于保存病例影像、用户头像和批改附件。本系统在配置中区分内部 \texttt{endpoint} 与公网 \texttt{public-endpoint}，后端上传后返回可访问 URL，避免医学附件直接占用数据库空间。

% ----------------------------------------------------------------------------
\section{容器化部署技术}
% 写作要点：
%   - Docker \cite{merkel2014docker}：一致的运行环境
%   - Docker Compose：多容器编排，本系统编排 mysql + redis + minio + backend + frontend

Docker 通过镜像封装应用及依赖，使开发、测试和部署环境保持一致\cite{merkel2014docker}。Docker Compose 用一份 \texttt{docker-compose.yml} 描述服务、端口、环境变量、网络和数据卷，适合本系统这种包含前端、后端、数据库、缓存和对象存储的组合部署。

本系统的 \texttt{docker-compose.yml} 编排 MySQL、Redis、MinIO、后端和前端五个服务。服务之间通过自定义网络通信，并通过 \texttt{depends\_on} 与 \texttt{healthcheck} 控制启动顺序。数据库密码、AI API Key、邮件服务密码等敏感参数通过 \texttt{.env} 注入，避免硬编码。

% ----------------------------------------------------------------------------
\section{本章小结}

本章介绍了系统所用的前端、后端、大语言模型、Agent、RAG、工具调用、数据存储和容器化部署技术，重点说明其在本系统中的作用。下一章将在这些技术基础上分析系统需求。
